%% ma: L12-B05-0035
%% lop: 12
%% chuong: 1
%% bai: 5
%% dang: 12.5.1
%% loai: TLN
%% muc: VD
%% dap_an: "1,2"
%% nguon: "HK1-DeOn_1_1.pdf (M10, Đề ôn tập học kì 1 Toán 12, Đề số 2), Phần III Câu 6"
%% kien_thuc: "Bài 2 (giá trị lớn nhất, nhỏ nhất), Bài 5 (ứng dụng đạo hàm giải quyết vấn đề thực tiễn)"
%% trang_thai: chinh-thuc
\begin{ex}
Một công ty sản xuất mỗi ngày được $x$ sản phẩm $(1\le x\le18)$. Tổng chi phí sản xuất $x$ sản phẩm tính bằng nghìn đồng cho bởi hàm chi phí $C(x)=x^3-3x^2+80x+500$. Giả sử công ty này bán hết sản phẩm mỗi ngày với giá $320$ nghìn đồng/sản phẩm. Gọi $B(x)$ là số tiền bán được và $L(x)$ là lợi nhuận thu được khi bán $x$ sản phẩm. Lợi nhuận (đơn vị tính triệu đồng) tối đa mà công ty có được là bao nhiêu?
\shortans{1,2}
\loigiai{$L(x)=B(x)-C(x)=320x-C(x)=-x^3+3x^2+240x-500$. $L'(x)=-3x^2+6x+240=0\Leftrightarrow x=10$ (loại $x=-8$).

$L(1)=-258$, $L(10)=1200$, $L(18)=-1040$. Lợi nhuận tối đa $1200$ nghìn đồng $=1{,}2$ triệu đồng.}
\end{ex}
